% GENERATED FILE - do not hand-edit.
% Source of truth: data/calculus-reviewed/c-10.json in the zenith-website repo.
% Regenerate with: python3 scripts/build_calculus_tex.py
\documentclass[11pt]{article}
\usepackage[utf8]{inputenc}
\usepackage[T1]{fontenc}
\usepackage{amsmath,amssymb}
\usepackage{graphicx}
\usepackage[margin=1in]{geometry}
\graphicspath{{../..//}}
\setlength{\parindent}{0pt}
\setlength{\parskip}{0.6em}
\begin{document}
\section*{Chapter 10: Sequences and Series --- Practice Exercises C10}

\bigskip\noindent\textbf{A1.}\ Which sequence converges?
\begin{list}{}{\setlength{\leftmargin}{3em}\setlength{\labelwidth}{2.5em}\setlength{\itemsep}{0pt}\setlength{\topsep}{2pt}}
\item[(A)] $a_n=n+3/n$
\item[(B)] $a_n=-1+(-1)^n/n$
\item[(C)] $a_n=\sin(n\pi/2)$
\item[(D)] $a_n=n!/3^n$
\end{list}

\bigskip\noindent\textbf{A2.}\ If $s_n=1+(-1)^n/n$, then
\begin{list}{}{\setlength{\leftmargin}{3em}\setlength{\labelwidth}{2.5em}\setlength{\itemsep}{0pt}\setlength{\topsep}{2pt}}
\item[(A)] $s_n$ diverges by oscillation
\item[(B)] $s_n$ converges to zero
\item[(C)] $\lim_{n\to\infty}s_n=1$
\item[(D)] $s_n$ diverges to infinity
\end{list}

\bigskip\noindent\textbf{A3.}\ The sequence $a_n = \sin\dfrac{n\pi}{6}$
\begin{list}{}{\setlength{\leftmargin}{3em}\setlength{\labelwidth}{2.5em}\setlength{\itemsep}{0pt}\setlength{\topsep}{2pt}}
\item[(A)] is unbounded
\item[(B)] converges to a number less than 1
\item[(C)] is bounded
\item[(D)] diverges to infinity
\end{list}

\bigskip\noindent\textbf{A4.}\ Which of the following sequences diverges?
\begin{list}{}{\setlength{\leftmargin}{3em}\setlength{\labelwidth}{2.5em}\setlength{\itemsep}{0pt}\setlength{\topsep}{2pt}}
\item[(A)] $a_n=(-1)^{n+1}/n$
\item[(B)] $a_n=2^n/e^n$
\item[(C)] $a_n=n^2/e^n$
\item[(D)] $a_n=n/\ln n$
\end{list}

\bigskip\noindent\textbf{A5.}\ The sequence $\{r^n\}$ converges if and only if
\begin{list}{}{\setlength{\leftmargin}{3em}\setlength{\labelwidth}{2.5em}\setlength{\itemsep}{0pt}\setlength{\topsep}{2pt}}
\item[(A)] $|r|<1$
\item[(B)] $|r|\leq 1$
\item[(C)] $-1<r\leq 1$
\item[(D)] $0<r<1$
\end{list}

\bigskip\noindent\textbf{A6.}\ $\displaystyle\sum u_n$ is a series of constants for which $\displaystyle\lim_{n\to\infty} u_n=0$. Which of the following statements is always true?
\begin{list}{}{\setlength{\leftmargin}{3em}\setlength{\labelwidth}{2.5em}\setlength{\itemsep}{0pt}\setlength{\topsep}{2pt}}
\item[(A)] $\displaystyle\sum u_n$ converges to a finite sum.
\item[(B)] $\displaystyle\sum u_n$ does not diverge to infinity.
\item[(C)] $\displaystyle\sum u_n$ is a positive series.
\item[(D)] None of statements (A), (B), or (C) are always true.
\end{list}

\bigskip\noindent\textbf{A7.}\ Using $\frac1{n(n+1)}=\frac1n-\frac1{n+1}$, evaluate $\sum_{n=1}^\infty\frac1{n(n+1)}$.
\begin{list}{}{\setlength{\leftmargin}{3em}\setlength{\labelwidth}{2.5em}\setlength{\itemsep}{0pt}\setlength{\topsep}{2pt}}
\item[(A)] $1$
\item[(B)] $3/2$
\item[(C)] $3/4$
\item[(D)] $\infty$
\end{list}

\bigskip\noindent\textbf{A8.}\ Find the sum of $2-1+\frac12-\frac14+\frac18-\cdots$.
\begin{list}{}{\setlength{\leftmargin}{3em}\setlength{\labelwidth}{2.5em}\setlength{\itemsep}{0pt}\setlength{\topsep}{2pt}}
\item[(A)] $4/3$
\item[(B)] $5/4$
\item[(C)] $1$
\item[(D)] $3/2$
\end{list}

\bigskip\noindent\textbf{A9.}\ Which of the following statements about series is true?
\begin{list}{}{\setlength{\leftmargin}{3em}\setlength{\labelwidth}{2.5em}\setlength{\itemsep}{0pt}\setlength{\topsep}{2pt}}
\item[(A)] If $\displaystyle\lim_{n\to\infty} u_n=0$, then $\displaystyle\sum u_n$ converges.
\item[(B)] If $\displaystyle\lim_{n\to\infty} u_n\neq 0$, then $\displaystyle\sum u_n$ diverges.
\item[(C)] If $\displaystyle\sum u_n$ diverges, then $\displaystyle\lim_{n\to\infty} u_n\neq 0$.
\item[(D)] $\displaystyle\sum u_n$ converges if and only if $\displaystyle\lim_{n\to\infty} u_n=0$.
\end{list}

\bigskip\noindent\textbf{A10.}\ Which of the following series diverges?
\begin{list}{}{\setlength{\leftmargin}{3em}\setlength{\labelwidth}{2.5em}\setlength{\itemsep}{0pt}\setlength{\topsep}{2pt}}
\item[(A)] $\sum1/n^2$
\item[(B)] $\sum\frac1{n^2+n}$
\item[(C)] $\sum\frac n{n^3+1}$
\item[(D)] $\sum\frac n{\sqrt{4n^2-1}}$
\end{list}

\bigskip\noindent\textbf{A11.}\ Which of the following series diverges?
\begin{list}{}{\setlength{\leftmargin}{3em}\setlength{\labelwidth}{2.5em}\setlength{\itemsep}{0pt}\setlength{\topsep}{2pt}}
\item[(A)] $3-1+1/9-1/27+\cdots$
\item[(B)] $1/\sqrt2-1/\sqrt3+1/\sqrt4-\cdots$
\item[(C)] $1/2^2-1/3^2+1/4^2-\cdots$
\item[(D)] $1-1.1+1.21-1.331+\cdots$
\end{list}

\bigskip\noindent\textbf{A12.}\ Let $S=\sum_{n=1}^\infty(2/3)^n$. Find $S$.
\begin{list}{}{\setlength{\leftmargin}{3em}\setlength{\labelwidth}{2.5em}\setlength{\itemsep}{0pt}\setlength{\topsep}{2pt}}
\item[(A)] $3/2$
\item[(B)] $4/3$
\item[(C)] $2$
\item[(D)] $3$
\end{list}

\bigskip\noindent\textbf{A13.}\ Which of the following expansions is impossible?
\begin{list}{}{\setlength{\leftmargin}{3em}\setlength{\labelwidth}{2.5em}\setlength{\itemsep}{0pt}\setlength{\topsep}{2pt}}
\item[(A)] $\sqrt{x-1}$ in powers of $x$
\item[(B)] $\sqrt{x+1}$ in powers of $x$
\item[(C)] $\ln x$ in powers of $(x-1)$
\item[(D)] $\tan x$ in powers of $(x-\pi/4)$
\end{list}

\bigskip\noindent\textbf{A14.}\ The series $\sum_{n=0}^\infty n!(x-3)^n$ converges if and only if
\begin{list}{}{\setlength{\leftmargin}{3em}\setlength{\labelwidth}{2.5em}\setlength{\itemsep}{0pt}\setlength{\topsep}{2pt}}
\item[(A)] $x=0$
\item[(B)] $2<x<4$
\item[(C)] $x=3$
\item[(D)] $2\leq x\leq4$
\end{list}

\bigskip\noindent\textbf{A15.}\ Let $f(x)=\displaystyle\sum_{n=0}^\infty x^n$. The radius of convergence of $\displaystyle\int_0^x f(t)\,dt$ is
\begin{list}{}{\setlength{\leftmargin}{3em}\setlength{\labelwidth}{2.5em}\setlength{\itemsep}{0pt}\setlength{\topsep}{2pt}}
\item[(A)] $0$
\item[(B)] $1$
\item[(C)] $2$
\item[(D)] $e$
\end{list}

\bigskip\noindent\textbf{A16.}\ The coefficient of $x^4$ in the Maclaurin series for $e^{-x/2}$ is
\begin{list}{}{\setlength{\leftmargin}{3em}\setlength{\labelwidth}{2.5em}\setlength{\itemsep}{0pt}\setlength{\topsep}{2pt}}
\item[(A)] $-1/24$
\item[(B)] $1/24$
\item[(C)] $-1/384$
\item[(D)] $1/384$
\end{list}

\bigskip\noindent\textbf{A17.}\ Evaluate $\int_0^{0.1}x^2e^{-x^2}\,dx$ to four decimal places using a series.
\begin{list}{}{\setlength{\leftmargin}{3em}\setlength{\labelwidth}{2.5em}\setlength{\itemsep}{0pt}\setlength{\topsep}{2pt}}
\item[(A)] $0.0002$
\item[(B)] $0.0003$
\item[(C)] $0.0032$
\item[(D)] $0.0033$
\end{list}

\bigskip\noindent\textbf{A18.}\ Use $\arctan1=1-1/3+1/5-1/7+\cdots$ to approximate $\pi/4$. What is the smallest number of terms guaranteed sufficient for error below0.001 by the first-omitted-term alternating-series bound?
\begin{list}{}{\setlength{\leftmargin}{3em}\setlength{\labelwidth}{2.5em}\setlength{\itemsep}{0pt}\setlength{\topsep}{2pt}}
\item[(A)] $200$
\item[(B)] $300$
\item[(C)] $400$
\item[(D)] $500$
\end{list}

\bigskip\noindent\textbf{A19.}\ Let $f(x)=\tan^{-1}x$ and $P_7(x)$ be the Taylor polynomial of degree 7 for $f$ about $x=0$. Given $\tan^{-1}x = \displaystyle\sum_{n=0}^\infty \dfrac{(-1)^n x^{2n+1}}{2n+1}$, it follows that if $-0.5<x<0.5$,
\begin{list}{}{\setlength{\leftmargin}{3em}\setlength{\labelwidth}{2.5em}\setlength{\itemsep}{0pt}\setlength{\topsep}{2pt}}
\item[(A)] $P_7(x)\leq \tan^{-1}x$
\item[(B)] $P_7(x)\geq \tan^{-1}x$
\item[(C)] $P_7(x)>\tan^{-1}x$ if $x<0$ but $P_7(x)<\tan^{-1}x$ if $x>0$
\item[(D)] $P_7(x)<\tan^{-1}x$ if $x<0$ but $P_7(x)>\tan^{-1}x$ if $x>0$
\end{list}

\bigskip\noindent\textbf{A20.}\ Let $f(x)=\tan^{-1}x$ and $P_{9}(x)$ be the Taylor polynomial of degree 9 for $f$ about $x=0$. Given $\tan^{-1}x = \displaystyle\sum_{n=0}^\infty \dfrac{(-1)^n x^{2n+1}}{2n+1}$, it follows that if $-0.5<x<0.5$,
\begin{list}{}{\setlength{\leftmargin}{3em}\setlength{\labelwidth}{2.5em}\setlength{\itemsep}{0pt}\setlength{\topsep}{2pt}}
\item[(A)] $P_{9}(x)\leq \tan^{-1}x$
\item[(B)] $P_{9}(x)\geq \tan^{-1}x$
\item[(C)] $P_{9}(x)>\tan^{-1}x$ if $x<0$ but $P_{9}(x)<\tan^{-1}x$ if $x>0$
\item[(D)] $P_{9}(x)<\tan^{-1}x$ if $x<0$ but $P_{9}(x)>\tan^{-1}x$ if $x>0$
\end{list}

\bigskip\noindent\textbf{A21.}\ Which of the following series converges?
\begin{list}{}{\setlength{\leftmargin}{3em}\setlength{\labelwidth}{2.5em}\setlength{\itemsep}{0pt}\setlength{\topsep}{2pt}}
\item[(A)] $\sum1/\sqrt[3]n$
\item[(B)] $\sum1/\sqrt n$
\item[(C)] $\sum1/(10n-1)$
\item[(D)] $\sum2/(n^2-5)$
\end{list}

\bigskip\noindent\textbf{A22.}\ Which of the following series diverges?
\begin{list}{}{\setlength{\leftmargin}{3em}\setlength{\labelwidth}{2.5em}\setlength{\itemsep}{0pt}\setlength{\topsep}{2pt}}
\item[(A)] $\sum_{n=1}^\infty\frac{n+1}{n!}$
\item[(B)] $\sum_{n=2}^\infty\frac1{n\ln n}$
\item[(C)] $\sum_{n=1}^\infty\frac{\ln n}{2^n}$
\item[(D)] $\sum_{n=1}^\infty\frac n{2^n}$
\end{list}

\bigskip\noindent\textbf{A23.}\ For which of the following series does the Ratio Test fail?
\begin{list}{}{\setlength{\leftmargin}{3em}\setlength{\labelwidth}{2.5em}\setlength{\itemsep}{0pt}\setlength{\topsep}{2pt}}
\item[(A)] $\sum1/n!$
\item[(B)] $\sum n/2^n$
\item[(C)] $1+1/2^{3/2}+1/3^{3/2}+1/4^{3/2}+\cdots$
\item[(D)] $\sum n^n/n!$
\end{list}

\bigskip\noindent\textbf{A24.}\ Which of the following alternating series diverges?
\begin{list}{}{\setlength{\leftmargin}{3em}\setlength{\labelwidth}{2.5em}\setlength{\itemsep}{0pt}\setlength{\topsep}{2pt}}
\item[(A)] $\sum\frac{(-1)^{n+1}(n-1)}{n+1}$
\item[(B)] $\sum\frac{(-1)^{n+1}}{\ln(n+1)}$
\item[(C)] $\sum\frac{(-1)^{n-1}}{\sqrt n}$
\item[(D)] $\sum\frac{(-1)^{n-1}n}{n^2+1}$
\end{list}

\bigskip\noindent\textbf{A25.}\ The power series $x+x^2/2+x^3/3+\cdots+x^n/n+\cdots$ converges if and only if
\begin{list}{}{\setlength{\leftmargin}{3em}\setlength{\labelwidth}{2.5em}\setlength{\itemsep}{0pt}\setlength{\topsep}{2pt}}
\item[(A)] $-1<x<1$
\item[(B)] $-1\leq x\leq 1$
\item[(C)] $-1\leq x<1$
\item[(D)] $-1<x\leq 1$
\end{list}

\bigskip\noindent\textbf{A26.}\ The series $(x+1)-(x+1)^2/2!+(x+1)^3/3!-(x+1)^4/4!+\cdots$ diverges
\begin{list}{}{\setlength{\leftmargin}{3em}\setlength{\labelwidth}{2.5em}\setlength{\itemsep}{0pt}\setlength{\topsep}{2pt}}
\item[(A)] for no real $x$
\item[(B)] if $-2<x\leq0$
\item[(C)] if $x<-2$ or $x>0$
\item[(D)] if $-2\leq x<0$
\end{list}

\bigskip\noindent\textbf{A27.}\ The series obtained by differentiating $(x-2)+(x-2)^2/4+(x-2)^3/9+(x-2)^4/16+\cdots$ term by term converges for
\begin{list}{}{\setlength{\leftmargin}{3em}\setlength{\labelwidth}{2.5em}\setlength{\itemsep}{0pt}\setlength{\topsep}{2pt}}
\item[(A)] $1\leq x\leq3$
\item[(B)] $1\leq x<3$
\item[(C)] $1<x\leq3$
\item[(D)] $0<x<4$
\end{list}

\bigskip\noindent\textbf{A28.}\ The degree3 Taylor polynomial at zero for $\sqrt{1+x}$ is
\begin{list}{}{\setlength{\leftmargin}{3em}\setlength{\labelwidth}{2.5em}\setlength{\itemsep}{0pt}\setlength{\topsep}{2pt}}
\item[(A)] $1+x/2-x^2/4+3x^3/8$
\item[(B)] $1+x/2-x^2/8+x^3/16$
\item[(C)] $1+x/2-x^2/8+x^3/8$
\item[(D)] $1-x/2+x^2/4-3x^3/8$
\end{list}

\bigskip\noindent\textbf{A29.}\ The Taylor polynomial of degree 3 at $x=1$ for $e^x$ is
\begin{list}{}{\setlength{\leftmargin}{3em}\setlength{\labelwidth}{2.5em}\setlength{\itemsep}{0pt}\setlength{\topsep}{2pt}}
\item[(A)] $e\left[1+(x-1)+\dfrac{(x-1)^2}{2}+\dfrac{(x-1)^3}{3}\right]$
\item[(B)] $e\left[1+(x+1)+\dfrac{(x+1)^2}{2!}+\dfrac{(x+1)^3}{3!}\right]$
\item[(C)] $e\left[1+(x-1)+\dfrac{(x-1)^2}{2!}+\dfrac{(x-1)^3}{3!}\right]$
\item[(D)] $e\left[1-(x-1)+\dfrac{(x-1)^2}{2!}+\dfrac{(x-1)^3}{3!}\right]$
\end{list}

\bigskip\noindent\textbf{A30.}\ The coefficient of $(x-\pi/4)^3$ in the Taylor series of $\cos x$ about $\pi/4$ is
\begin{list}{}{\setlength{\leftmargin}{3em}\setlength{\labelwidth}{2.5em}\setlength{\itemsep}{0pt}\setlength{\topsep}{2pt}}
\item[(A)] $-1/12$
\item[(B)] $1/12$
\item[(C)] $1/(6\sqrt2)$
\item[(D)] $-1/(3\sqrt2)$
\end{list}

\bigskip\noindent\textbf{A31.}\ The coefficient of $x^2$ in the Maclaurin series for $e^{\sin x}$ is
\begin{list}{}{\setlength{\leftmargin}{3em}\setlength{\labelwidth}{2.5em}\setlength{\itemsep}{0pt}\setlength{\topsep}{2pt}}
\item[(A)] $0$
\item[(B)] $1$
\item[(C)] $1/2!$
\item[(D)] $1/4$
\end{list}

\bigskip\noindent\textbf{A32.}\ The coefficient of $(x-1)^5$ in the Taylor series for $x\ln x$ about1 is
\begin{list}{}{\setlength{\leftmargin}{3em}\setlength{\labelwidth}{2.5em}\setlength{\itemsep}{0pt}\setlength{\topsep}{2pt}}
\item[(A)] $-1/20$
\item[(B)] $-1/5!$
\item[(C)] $-1/4!$
\item[(D)] $-1/5$
\end{list}

\bigskip\noindent\textbf{A33.}\ The radius of convergence of $\sum_{n=1}^\infty\frac{x^n}{2^n}\frac{n^n}{n!}$ is
\begin{list}{}{\setlength{\leftmargin}{3em}\setlength{\labelwidth}{2.5em}\setlength{\itemsep}{0pt}\setlength{\topsep}{2pt}}
\item[(A)] $2$
\item[(B)] $2/e$
\item[(C)] $e/2$
\item[(D)] $\infty$
\end{list}

\bigskip\noindent\textbf{A34.}\ The Taylor polynomial of degree 3 at $x=0$ for $(1+x)^p$, where $p$ is a constant, is
\begin{list}{}{\setlength{\leftmargin}{3em}\setlength{\labelwidth}{2.5em}\setlength{\itemsep}{0pt}\setlength{\topsep}{2pt}}
\item[(A)] $1+px+p(p-1)x^2+p(p-1)(p-2)x^3$
\item[(B)] $1+px+\dfrac{p(p-1)}{2}x^2+\dfrac{p(p-1)(p-2)}{3}x^3$
\item[(C)] $1+px+\dfrac{p(p-1)}{2!}x^2+\dfrac{p(p-1)(p-2)}{3!}x^3$
\item[(D)] $px+\dfrac{p(p-1)}{2!}x^2+\dfrac{p(p-1)(p-2)}{3!}x^3$
\end{list}

\bigskip\noindent\textbf{A35.}\ The Taylor series for $\ln(1+2x)$ about zero is
\begin{list}{}{\setlength{\leftmargin}{3em}\setlength{\labelwidth}{2.5em}\setlength{\itemsep}{0pt}\setlength{\topsep}{2pt}}
\item[(A)] $2x-2x^2+8x^3/3-4x^4+\cdots$
\item[(B)] $2x-4x^2+16x^3+\cdots$
\item[(C)] $2x-x^2+8x^3/3-4x^4+\cdots$
\item[(D)] $2x-2x^2+4x^3/3-2x^4/3+\cdots$
\end{list}

\bigskip\noindent\textbf{A36.}\ For which values of $x$ does $\sum_{n=1}^\infty n2^n/x^n$ converge?
\begin{list}{}{\setlength{\leftmargin}{3em}\setlength{\labelwidth}{2.5em}\setlength{\itemsep}{0pt}\setlength{\topsep}{2pt}}
\item[(A)] only $x=0$
\item[(B)] $-2<x<2$
\item[(C)] $|x|>2$
\item[(D)] $|x|\geq2$
\end{list}

\bigskip\noindent\textbf{A37.}\ The third-degree Taylor polynomial $P_3(x)$ for $\sin x$ about $\dfrac{\pi}{4}$ is
\begin{list}{}{\setlength{\leftmargin}{3em}\setlength{\labelwidth}{2.5em}\setlength{\itemsep}{0pt}\setlength{\topsep}{2pt}}
\item[(A)] $\dfrac{1}{\sqrt{2}}\left(\left(x-\dfrac{\pi}{4}\right)-\dfrac{1}{3!}\left(x-\dfrac{\pi}{4}\right)^3\right)$
\item[(B)] $\dfrac{1}{\sqrt{2}}\left(1+\left(x-\dfrac{\pi}{4}\right)-\dfrac{1}{2}\left(x-\dfrac{\pi}{4}\right)^2+\dfrac{1}{3!}\left(x-\dfrac{\pi}{4}\right)^3\right)$
\item[(C)] $\dfrac{1}{\sqrt{2}}\left(1+\left(x-\dfrac{\pi}{4}\right)-\dfrac{1}{2!}\left(x-\dfrac{\pi}{4}\right)^2-\dfrac{1}{3!}\left(x-\dfrac{\pi}{4}\right)^3\right)$
\item[(D)] $1+\left(x-\dfrac{\pi}{4}\right)-\dfrac{1}{2}\left(x-\dfrac{\pi}{4}\right)^2-\dfrac{1}{6}\left(x-\dfrac{\pi}{4}\right)^3$
\end{list}

\bigskip\noindent\textbf{A38.}\ If $h(1)=h\prime(1)=h\prime\prime(1)=h\prime\prime\prime(1)=6$, which cubic is its degree3 Taylor polynomial about1?
\begin{list}{}{\setlength{\leftmargin}{3em}\setlength{\labelwidth}{2.5em}\setlength{\itemsep}{0pt}\setlength{\topsep}{2pt}}
\item[(A)] $6+6x+6x^2+6x^3$
\item[(B)] $6+6(x-1)+6(x-1)^2+6(x-1)^3$
\item[(C)] $6+6x+3x^2+x^3$
\item[(D)] $6+6(x-1)+3(x-1)^2+(x-1)^3$
\end{list}

\bigskip\noindent\textbf{B1.}\ Which of the following statements about series is false?
\begin{list}{}{\setlength{\leftmargin}{3em}\setlength{\labelwidth}{2.5em}\setlength{\itemsep}{0pt}\setlength{\topsep}{2pt}}
\item[(A)] $\displaystyle\sum_{k=1}^\infty u_k = \displaystyle\sum_{k=m}^\infty u_k$, where $m$ is any positive integer.
\item[(B)] If $\displaystyle\sum u_n$ converges, so does $\displaystyle\sum cu_n$ if $c\neq 0$.
\item[(C)] If $\displaystyle\sum a_n$ and $\displaystyle\sum b_n$ converge, so does $\displaystyle\sum(ca_n+b_n)$, where $c\neq 0$.
\item[(D)] Rearranging the terms of a positive convergent series will not affect its convergence or its sum.
\end{list}

\bigskip\noindent\textbf{B2.}\ Which of the following statements is always true?
\begin{list}{}{\setlength{\leftmargin}{3em}\setlength{\labelwidth}{2.5em}\setlength{\itemsep}{0pt}\setlength{\topsep}{2pt}}
\item[(A)] If $\displaystyle\sum u_n$ converges, then so does the series $\displaystyle\sum |u_n|$.
\item[(B)] If a series is truncated after the $n$th term, then the error is less than the first term omitted.
\item[(C)] If the terms of an alternating series decrease, then the series converges.
\item[(D)] None of statements (A), (B), or (C) are always true.
\end{list}

\bigskip\noindent\textbf{B3.}\ Which of the following series can be used to compute $\ln 0.8$?
\begin{list}{}{\setlength{\leftmargin}{3em}\setlength{\labelwidth}{2.5em}\setlength{\itemsep}{0pt}\setlength{\topsep}{2pt}}
\item[(A)] $\ln(x-1)$ expanded about $x=0$
\item[(B)] $\ln x$ about $x=0$
\item[(C)] $\ln x$ expanded about $x=1$
\item[(D)] $\ln(x-1)$ expanded about $x=1$
\end{list}

\bigskip\noindent\textbf{B4.}\ Let $f(x)=\displaystyle\sum_{n=0}^\infty a_nx^n$, $g(x)=\displaystyle\sum_{n=0}^\infty b_nx^n$. Suppose both series converge for $|x|<R$. Let $x_0$ be a number such that $|x_0|<R$. Which of statements (A)--(C) is false?
\begin{list}{}{\setlength{\leftmargin}{3em}\setlength{\labelwidth}{2.5em}\setlength{\itemsep}{0pt}\setlength{\topsep}{2pt}}
\item[(A)] $\displaystyle\sum_{n=0}^\infty (a_n+b_n)(x_0)^n$ converges to $f(x_0)+g(x_0)$.
\item[(B)] $f(x)=\displaystyle\sum_{n=0}^\infty a_nx^n$ is continuous at $x=x_0$.
\item[(C)] $\displaystyle\sum_{n=1}^\infty na_nx_0^{n-1}$ converges to $f'(x_0)$.
\item[(D)] Statements (A)--(C) are all true.
\end{list}

\bigskip\noindent\textbf{B5.}\ If $\sin x$ is approximated by $x-x^3/3!$ for $|x|<1$, then the error is numerically less than
\begin{list}{}{\setlength{\leftmargin}{3em}\setlength{\labelwidth}{2.5em}\setlength{\itemsep}{0pt}\setlength{\topsep}{2pt}}
\item[(A)] $0.003$
\item[(B)] $0.005$
\item[(C)] $0.008$
\item[(D)] $0.009$
\end{list}

\bigskip\noindent\textbf{B6.}\ Let $f(x)=\sum_{n=0}^\infty a_nx^n$ satisfy $f\prime(x)=-f(x)$ and $f(0)=1$. How many terms are needed to find $f(0.2)$ correct to three decimal places?
\begin{list}{}{\setlength{\leftmargin}{3em}\setlength{\labelwidth}{2.5em}\setlength{\itemsep}{0pt}\setlength{\topsep}{2pt}}
\item[(A)] $2$
\item[(B)] $3$
\item[(C)] $4$
\item[(D)] $5$
\end{list}

\bigskip\noindent\textbf{B7.}\ The series $\sum_{n=1}^\infty(\pi^3/(3\pi))^n$
\begin{list}{}{\setlength{\leftmargin}{3em}\setlength{\labelwidth}{2.5em}\setlength{\itemsep}{0pt}\setlength{\topsep}{2pt}}
\item[(A)] equals1
\item[(B)] equals $3\pi/(\pi^3-3\pi)$
\item[(C)] equals $\pi^3/(3\pi-\pi^3)$
\item[(D)] diverges
\end{list}

\bigskip\noindent\textbf{B8.}\ When $\sum_{n=1}^\infty(-1)^{n-1}/(3n-1)$ is approximated by its first300 terms, the error is closest to
\begin{list}{}{\setlength{\leftmargin}{3em}\setlength{\labelwidth}{2.5em}\setlength{\itemsep}{0pt}\setlength{\topsep}{2pt}}
\item[(A)] $0.001$
\item[(B)] $0.002$
\item[(C)] $0.005$
\item[(D)] $0.01$
\end{list}

\bigskip\noindent\textbf{B9.}\ To estimate $e^x$ uniformly for $|x|\leq2$ with error below0.001, the Lagrange bound suggests a Taylor polynomial at zero of degree at least
\begin{list}{}{\setlength{\leftmargin}{3em}\setlength{\labelwidth}{2.5em}\setlength{\itemsep}{0pt}\setlength{\topsep}{2pt}}
\item[(A)] $9$
\item[(B)] $10$
\item[(C)] $11$
\item[(D)] $12$
\end{list}

\end{document}
